\section*{After the Mass: The Leonine Prayers}
\addcontentsline{toc}{section}{After the Mass: The Leonine Prayers}

\begin{massbox}{Boundary before interpretation}
The Leonine prayers are not a concluding movement of Mass. In its ordinary printed sequence the 1962 typical Missal's \latin{Ordo Missae} ends on page~328 with the Last Gospel, normally John~1:1--14, and \latin{R. Deo gratias}; General Rubric 510 nevertheless omits the Last Gospel in several listed cases. The prayers below were imposed by external law \emph{after} Masses within their scope. An 1887 response required them to follow immediately when the Last Gospel had been completed, but that wording must not be converted into proof that the annex applied identically to every 1962 Last-Gospel exception. They were historically adjacent and obligatory within their juridical scope without becoming part of the Eucharistic sacrifice or the Missal's Order.
\end{massbox}

\subsection*{The received post-Mass sequence}

The form commonly encountered at the 1962 horizon began with three \latin{Ave Maria}, continued with the \latin{Salve Regina}, its versicle \latin{Ora pro nobis, sancta Dei Genetrix} and response, and a collect beginning \latin{Deus, refugium nostrum et virtus}. It then included the short prayer \latin{Sancte Michael Archangele, defende nos in proelio}, asking the archangel's defense against the devil's wickedness and snares. The invocation \latin{Cor Iesu Sacratissimum, miserere nobis} was also customarily appended three times under the discipline explained below.

This outline identifies the received sequence without pretending that every element had the same juridical origin. The three Aves, Marian antiphon, versicle, and collect belonged to Leo~XIII's 1884 universal prescription. The revised collect and short Saint Michael prayer came with the revised form ordered in 1886. Pius~X desired uniform use of the Sacred Heart invocation and encouraged priests to say it, but the controlling declaration expressly denied that he had imposed an obligation in the proper sense. Local legislation and customary manuals could therefore present a uniform practical text whose components still possessed different legal histories.

The prayer's ordinary bodily grammar was kneeling. This posture did not extend the Canon's sacrificial action or add another liturgical climax. It marked supplication after the dismissal: the Church that had received blessing remained conscious of conflict, sin, persecution, and the need for protection as she returned to the world.

\subsection*{From a threatened temporal order to universal supplication}

The Sacred Congregation of Rites decree \emph{Iam inde ab anno}, dated 6~January 1884, looked back to an 1859 prescription of Pius~IX for churches in the Papal territory during grave circumstances. That retrospective report is secure evidence for the antecedent; it is not a warrant to reconstruct the exact 1859 formulary without the original act.

Leo~XIII universalized and specified the discipline. The decree ordered that, \latin{in fine cuiusque Missae sine cantu celebratae, flexis genibus recitentur}---at the end of every Mass celebrated without chant, the prayers were to be recited kneeling (\emph{Acta Sanctae Sedis} 16 [1883--84], 239--240). The exact legal phrase is “Mass celebrated without chant.” Calling the observance the prayers after Low Mass is convenient and historically familiar, but it must not replace the decree's wording or erase later interpretive exceptions.

The 1884 sequence comprised three Hail Marys, the \latin{Salve Regina}, versicle and response, and a collect invoking the intercession of Mary, Joseph, Peter, Paul, and all saints for the Church's needs. In 1886 Leo ordered a revised form: the collect explicitly sought the conversion of sinners and the liberty and exaltation of holy Mother Church, and the short Saint Michael prayer was added. A contemporary official diocesan circular securely witnesses reception of that new form. The central Roman promulgation locus for the 1886 wording has not been recovered in this audit; that limit counsels precision about the fact of revision without manufacturing a citation.

An 1887 response of the Sacred Congregation of Rites supplies an especially important boundary. Asked whether another rite could intervene, it answered that the prayers were to be recited \latin{immediate expleto ultimo Evangelio}---immediately when the Last Gospel had been completed (\emph{Acta Sanctae Sedis} 23 [1890--91], 128--129). “Completed” and “immediately” belong together: the Mass had ended, and the annex was nevertheless not to be casually displaced. The response states the normal Last-Gospel branch; it does not by itself resolve every later Mass in which General Rubric 510 omits that Gospel.

\subsection*{Mary, Michael, and the Sacred Heart}

The Marian sequence is not a rival epilogue to Christ's sacrifice. Gabriel's greeting and Elizabeth's blessing in Luke~1:26--45 confess what grace has done in Mary and name the fruit of her womb; the \latin{Salve Regina} addresses the Mother of Mercy as an intercessor who shows exiles Jesus, the blessed fruit. Her mediation is maternal, participatory, and wholly subordinate to the one mediation of Christ (1~Tim.~2:5; \emph{Lumen gentium} 60--62). Coming after Communion and mission, the prayer asks the foremost member of the redeemed Church to aid pilgrims in living from the Son they have received.

Saint Michael's biblical setting is conflict under divine sovereignty. Daniel~10 and 12 present him as a princely defender of God's people; Jude~9 shows him opposing the devil without autonomous judgment; Revelation~12:7--12 sets angelic warfare inside the victory of the Lamb and the martyrs' testimony. The prayer therefore asks angelic defense, not a second redemption and not a magical contest between nearly equal powers. Aquinas places angelic ministry beneath God's providence: angels truly act as intellectual creatures and ministering spirits, while every power by which they defend remains created and governed by God (\emph{ST} I, qq.~110--114).

On 17~June 1904 an act of the Sacred Congregation of Indulgences and Sacred Relics permitted the invocation \latin{Cor Iesu Sacratissimum, miserere nobis} to be added three times to the prayers and attached an indulgence to its threefold recitation (\emph{Acta Sanctae Sedis} 36 [1903--04], 750). The same dicastery's declaration of 19~August stated that an obligation in the proper sense had not been imposed by Pius~X; he desired uniformity and priests were to be encouraged to recite it (\emph{Acta Sanctae Sedis} 37 [1904--05], 125--126). It is therefore inaccurate to flatten an indulgenced and strongly encouraged threefold practice, a local command, and universal precept into the assertion that Pius~X universally made the invocation obligatory. The Sacred Heart turns the final plea to the incarnate charity symbolized by the pierced side (John~19:34--37): spiritual combat is governed by the love of the Lamb, not fascination with demonic power.

\subsection*{The intention for Russia}

On 30~June 1930 Pius~XI's allocution \emph{Indictam ante} redirected the existing prayers to a named emergency. He asked that Christ restore tranquility and the freedom to profess the faith to the afflicted children of Russia, and directed Latin priests throughout the world to say the prayers after Mass \latin{ad hanc ipsam mentem, scilicet pro Russia}---for this intention, namely for Russia (AAS 22 [1930], 301). On 11~July the Pontifical Commission for Russia exhorted the non-Latin priests under its authority to commend the same intention during the Divine Liturgy and prescribed additions to the Great Litany and the prayer after the ambo (AAS 22 [1930], 366). That was an analogous Eastern implementation, not the source of the already worldwide Latin direction.

Later speech often compresses this history into “for the conversion of Russia.” Conversion is already present in the 1886 collect's prayer for sinners, but it is not a substitute quotation for Pius~XI's more concrete 1930 language about tranquility and freedom to profess the faith. Nor does the official act say that the decision was made in documentary response to Fatima. The theological relation may be contemplated only after the historical claims have been kept distinct; the promulgated reason and intention must be allowed to speak in their own terms.

This redirection also reveals the prayers' catholic reach. A short devotion after local celebrations made the persecution of distant Christians an object of persevering common petition. First Timothy~2:1--6 places intercession for public order beneath God's will to save and Christ's unique mediation. The prayers did not claim a sacramental change in Russia; they exercised the Church's real but dependent agency by asking omnipotent providence for conversion, liberty, endurance, and peace.

\subsection*{What was true in 1962}

Three propositions must be held together. First, the typical Missal's \latin{Ordo Missae} did not print the Leonine prayers and ordinarily ended with the Last Gospel. Second, the external universal prescription had not yet been suppressed, so the prayers remained part of the discipline surrounding Masses within its scope in 1962. Third, the Sacred Congregation of Rites' decree of 9~March 1960, approved by John~XXIII, supplied controlling exceptions current at that horizon (AAS 52 [1960], 360).

That decree permitted omission after a Nuptial Mass or a Mass for First Communion, General Communion, Confirmation, Ordination, or religious profession; when another function or pious exercise immediately and duly followed Mass; when a homily occurred within Mass; and at a dialogue Mass on Sundays and feast days. It also allowed Ordinaries to approve vernacular texts. These exceptions belong to the annex's own law, not to General Rubric 510's independent list of Masses omitting the Last Gospel. Thus neither “every Low Mass” nor “every completed Last Gospel” is an adequate 1962 rule by itself.

This distinction between text and law is ordinary Catholic reasoning. A Missal governs the rite it prints; competent authority can also attach an observance before or after that rite and specify its exceptions. The attached act may be obligatory without becoming sacramental matter, Eucharistic prayer, or an intrinsic condition of a valid Mass. Conversely, absence from the printed \latin{Ordo} did not make the 1962 obligation imaginary. Edition history and juridical history intersect without becoming identical.

\subsection*{Suppression: promulgation and effective date}

The instruction \emph{Inter Oecumenici}, promulgated 26~September 1964, stated at no.~48(j): \latin{Ultimum Evangelium omittitur; preces Leonianae supprimuntur}---the Last Gospel is omitted; the Leonine prayers are suppressed (AAS 56 [1964], 888). The same instruction's concluding norm made its provisions effective on 7~March 1965, the First Sunday of Lent. The accurate chronology is therefore promulgation in 1964 and operative suppression in 1965.

Suppression ended the juridical prescription; it did not condemn the Hail Mary, \latin{Salve Regina}, Saint Michael prayer, Sacred Heart devotion, or prayer for persecuted Christians. Neither does the former obligation determine every present disciplinary question. Voluntary devotional use after Mass, where otherwise lawful and pastorally ordered, is a question distinct from claiming that the suppressed universal rule still binds.

\subsection*{Two historical safeguards}

First, no contemporary primary evidence located in this audit supports the popular story that Leo~XIII composed the short Saint Michael prayer after hearing a post-Mass dialogue between God and Satan or receiving a precisely dated vision. Versions of the story conflict about date and content and appear in much later testimony. The prayer's sound doctrine does not need a legend to authenticate it, and historical piety is served by truth rather than an edifying invention.

Second, the short \latin{Sancte Michael Archangele} attached to the Leonine sequence must not be confused with the much longer \latin{Exorcismus in satanam et angelos apostaticos} published under Leo~XIII in 1890 (\emph{Acta Sanctae Sedis} 23 [1890--91], 743--746). The former is a public supplication for defense after Mass; the latter belongs to a distinct juridical and ritual context. Shared imagery and papal era do not make the texts or permissions interchangeable.

\begin{studybox}{Relation to the whole Mass}
The Leonine sequence asks that Eucharistic mission be guarded in history: Mary embodies redeemed receptivity and maternal intercession; Michael signifies created strength obedient to providence; the Sacred Heart names the incarnate charity from which victory flows; the Russian intention refuses indifference to a persecuted body of Christ. Yet every petition depends on what the Mass has already made present. The prayers add no victim, priesthood, consecration, or Communion. Wherever the external prescription applied, they began only after the Mass's final printed action---ordinarily the Last Gospel and \latin{Deo gratias}---when the Church, blessed and sent, confessed that mission remained dependent on grace.
\end{studybox}
